\documentclass[12pt, a4paper]{article}
\usepackage[T1]{fontenc}
\usepackage{lmodern}
\usepackage{amsmath, amsthm, amssymb, mathtools}
\usepackage[margin=1in]{geometry}
\usepackage{xr-hyper}
\usepackage[colorlinks=true,linkcolor=blue,citecolor=blue]{hyperref}
\usepackage{cleveref}
\usepackage{graphicx, subcaption}
\usepackage{natbib}
\usepackage{booktabs}
\usepackage{enumerate}

% Code references: scaled monospace with line-breaking
\newcommand{\code}[1]{\texttt{\small #1}}

\newtheorem{theorem}{Theorem}[section]
\newtheorem{lemma}[theorem]{Lemma}
\newtheorem{proposition}[theorem]{Proposition}
\newtheorem{corollary}[theorem]{Corollary}
\newtheorem{definition}{Definition}[section]
\newtheorem{remark}{Remark}[section]
\newtheorem{conjecture}[theorem]{Conjecture}
\newtheorem{derivation}[theorem]{Derivation}
\newtheorem{observation}[theorem]{Observation}
\numberwithin{equation}{section}

\newcommand{\dd}{\mathrm{d}}
\newcommand{\seriestitle}{Why Rotating Fluids Make Polygons}


\externaldocument[I-]{../paper-1-mathematics/main}
\externaldocument[II-]{../paper-2-physics/main}
\externaldocument[III-]{../paper-3-planets/main}
\externaldocument[A-]{../paper-A-appendices/main}

\title{\seriestitle\\[6pt]\large Overview and Preface}
\author{Gordon Speagle}
\date{}

\begin{document}
\maketitle

\begin{abstract}
On the flat ($K=0$) plane, the regular $N$-gon of point vortices is an
energy minimum on the surface of conserved angular impulse
$L = \sum|z_k|^2$ and linear impulse $\mathbf{P} = \sum z_k = 0$ for
all $N \le 7$, and unstable for $N \ge 8$ without a central vortex.
This stability boundary is an invariant of
the two-dimensional Laplacian: the logarithmic Green's function
$G \propto -\ln r$ is the unique radially symmetric planar interaction
for which the Havelock eigenvalue $\lambda_m = (N-1) - m(N-m)/2$ is
scale-independent, and $N = 7$ is the threshold where the leading
eigenvalue $\lambda_3 = 0$ (the linear restoring force vanishes).
A quartic calculation on the constraint surface confirms that the
$N = 7$ ring is nonlinearly stable: the quartic coefficient
$\alpha_0 = 45/14 \approx 3.21 > 0$ (in the normal form
$\mathcal{H} = (\alpha_0/4)\varepsilon^4$) supplies a restoring force
where the linear one is absent.

For $N \ge 8$ the stability boundary requires a central vortex of
strength $\kappa_0/\kappa \ge \kappa_{\mathrm{crit}}(N)$.  The sharpest
prediction is at $N = 8$: $\kappa_{\mathrm{crit}}(8) = 1/2$.  In a
Bose--Einstein condensate, where circulation is quantized, this is an
integer-valued threshold: a ring of eight unit-winding vortices is
stable if and only if the central vortex carries winding number
$q_0 \ge 1$.

The central unifying result is a \emph{constraint-intersection
principle}: three selection mechanisms constrain the polygon
wavenumber $N$ from above, and the observed $N$ is determined by
whichever is tightest.
Two are fully derived from the 2D Laplacian~$\nabla^2$:
\emph{Thomson stability} (eigenvalues of the Green's function
$G(\nabla^2)$) and \emph{Rossby stationarity} (stationary solutions
of $(\nabla^2 + \beta\,\partial_y)\psi = 0$).
The third---\emph{geometric packing}---has a Laplacian-derived
skeleton (the domain admissibility condition
$\sin(\pi/N) \ge \varepsilon/R$ for the Green's function on the
punctured domain) but requires one empirical input: the
$\beta$-drift enlargement $r_{\mathrm{excl}} \approx
1.3\,r_{\mathrm{core}}$ from nonlinear vortex dynamics on a
rotating sphere \citep{Gavriel2021}, not derivable from the
linear operator alone.
Saturn's hexagon ($N = 6$) is Rossby-limited (fully derived).
Jupiter's north-pole octagon ($N = 8$) is Thomson-limited
(fully derived; the prediction $N_{\max} = 8$ relies on
$\kappa_0/\kappa \lesssim 1$, testable with improved Juno
wind measurements).
Jupiter's south-pole pentagon ($N = 5$) is packing-limited
(Laplacian skeleton plus the empirical $1.3$ factor).

On the hyperbolic plane $\mathbf{H}^2$ with curvature $K = -1/a^2$,
the stability boundary rises monotonically: $N_{\mathrm{crit}} = 12$
at ring-to-curvature ratio $R/a = 1$---a falsifiable prediction
for circuit-QED hyperbolic lattices \citep{Kollar2019}.
The $N \to (N{+}1)$ transition threshold $\xi^*(N)$ satisfies
a \emph{palindromic} quadratic $A\xi^2 + B\xi + A = 0$, whose
equal leading and trailing coefficients force
$\xi \cdot \xi' = 1$---the roots are multiplicative inverses.
This is not a computational coincidence: it encodes the conformal
inversion $\xi \leftrightarrow 1/\xi$ of the Poincar\'e disk
(Paper~I, Proposition~2.3), connecting stability thresholds
to the Pell equation $a^2 - Db^2 = 1$ of algebraic number theory.

The threshold $\xi^*(N)$ is an algebraic integer---living in the
maximal order of $\mathbb{Q}(\sqrt{D})$---for exactly five values:
$N \in \{8, 9, 11, 15, 23\}$ (characterised by $(N{-}7) \mid 16$;
differences $\{1,2,4,8\}$, powers of~$2$).  At each, $\xi^*$ is
a power of the inverse fundamental unit: $\xi^*(8) = (8{+}3\sqrt{7})^{-1}$
in $\mathbb{Z}[\sqrt{7}]$, $\xi^*(23) = \varphi^{-2}$ where
$\varphi = (1{+}\sqrt{5})/2$ is the golden ratio.
The monic polynomial traces $\{16, 10, 6, 4, 3\}$ follow a binary
subdivision $c(k) = 2 + 2^{4-k}$ at $N = 7 + 2^k$, converging to
the golden ratio ($c = 3$) as the algebraically simplest
palindromic unit.  Each threshold has continued fraction
$1/\xi^* = [c{-}1;\;\overline{1,\, c{-}2}]$, degenerating to
\end{abstract}

\section*{Preface}

This project began years ago when I first encountered the Voyager
photographs of Saturn's polar hexagonal storm.  Why a polygonal ring?
Why the persistence?  I researched the relevant evidence and arguments,
and while informative, the underlying framework did not adequately
satiate my desire for a clear answer.  I catalogued it as one of
those interesting intersections between the physical world and
geometry---crystals, viruses, snowflakes, coral reefs, beehives,
nautilus shells---many of which have been rigorously explained, yet
the explanatory framework for polygonal formation in the atmospheres
of the two largest planets in our Solar System remained disconnected
and incomplete.

My educational background is primarily philosophy with a mathematical
timbre.  Formally, I am more familiar with Peano, Frege, Russell, and
G\"odel than the applied disciplines.  A hobbyist's obsession with
number theory is an apt description for my post-educational years.
Professionally, I am a full-stack developer working in healthcare.

Fast-forwarding: I was watching an astronomy video with my children and
the polar polygons on Jupiter and Saturn immediately astonished them
and once again piqued my curiosity.  I dove back into the most current
research.  The Cassini and Juno images hinted at a deeper mathematical
connection between the multiple $N$-gons on the gas giants.  The hook
was set.  I ordered a copy of \emph{Geophysical Fluid Dynamics} by
Joseph Pedlosky and started writing code.

Observationally, the storms show central vortices and logarithmic
spirals eventually mutating into polygonal rings.  From my wanderings
in the mathematical literature, I knew there was a formal connection
between the first two via M\"obius transformations.  This was my
mathematical starting point and eventual foundation for the resulting
paper.

I became determined to pull the connecting string until the knot
unravelled.  This project has crossed many mathematical, physical,
and astrophysical domains as a result.  I am by no means a subject
matter expert in all of these fields; my research directions were
motivated by the singular goal of finding an answer to the original
question, wherever that led.  Institutional research correctly forces
disciplinary depth; the cost is that cross-domain connections can go
unnoticed.  I am unaffiliated, which meant I could follow the
mathematics wherever it led, without the professional incentive to
stay within a single field.  Whether that freedom produced something
worth the lack of expert supervision is for the audience to assess.

My process was iterative: at each step I identify a potential connecting
point, honestly assess what I do and do not understand about the
underlying structures, research the applicable domain, study that
specific subset until I have an intuitive understanding in the context
of my project, then develop and test software that confirms my
understanding and verifies against existing mathematics, physics, or
astrophysics.  At that point I architect code that will support, prove,
or disprove my connective hypotheses.  Rinse and repeat.

My background gave me a solid logical foundation of what constitutes
a ``proof'' in the strictest sense.  This mapped well onto the sections
of the paper that required such reasoning.  The software allowed me to
quickly assess feasibility, routes to proving, potential logical errors
or contradictions, and blind spots.

Full transparency: the review process drew on AI-assisted critique.
With Claude Code, I built structured review agents---a professional
mathematician, a physicist, and a technical writer.  I needed my
subjective understanding, writing, and results to be in a language
conversant with the modern research community.  Once I had my ideas
fully fleshed out in writing, I used the agents to scrutinise, review,
and find logical errors.  The philosophical ideas, logic, proof
scaffolding, development, and insights from the paper are my own.
This process repeated until I was confident that the paper accurately
reflected the quality, tone, and accuracy expected of a published
article in academic journals.

I have made a serious attempt to cite and attribute assiduously, as
well as to honestly assess any weakness in argumentation.
Acknowledging vulnerabilities, although in a sense ``knee-capping''
portions of the paper, is intellectually honest, and since my goal was
to answer \emph{why}, I do not think the acknowledged shortcomings
detract from the overall aim.

The paper is admittedly dense in places and broad in scope.  No part
requires fluency in all parts.  The mathematician will be less inclined
to devote attention to planetary atmospheric dynamics.  The planetary
scientist might find the algebraic number field portions less
interesting than the mechanisms at play in polar storms.
Experimentalists in circuit QED will be drawn to the BEC predictions.
I have attempted to build the paper logically, starting with the
mathematics first and progressing to the applied planetary physics.
The table of contents will direct the reader to the areas most
relevant to their interests.

Feedback, comments, and suggestions are welcome.

\begin{flushright}
\emph{Gordon Speagle}
\end{flushright}

\section{Guide to the Series}

This work is presented as a series of four papers plus appendices:
\begin{description}
\item[Paper~I: Mathematical Foundations.] The logarithmic interaction,
  sign rule, Havelock identity, stability thresholds on curved surfaces,
  and algebraic structure (Pell equations, golden ratio, continued fractions).
\item[Paper~II: Classical and Quantum Physics.] Point vortex dynamics,
  constrained energy landscape, central vortex stabilization, $N=7$
  bifurcation, multi-ring dichotomy, blob bridge, $K_0$ interaction,
  disorder resilience, and BEC predictions.
\item[Paper~III: Planetary and Astrophysical Applications.] The
  constraint-intersection framework, Saturn hexagon, Jupiter polygons,
  Neptune predictions, first-principles prediction chain, and the
  Onsager--Arnold two-timescale resolution.
\item[Appendices.] Technical derivations, code availability, and
  open problems for future research.
\end{description}

Each paper can be read independently; cross-references to results
in other papers are marked with the paper number (e.g., Theorem~I-2.1
refers to Theorem~2.1 in Paper~I).

\bibliographystyle{plainnat}
\bibliography{../shared/refs}
\end{document}
